% !TEX program = pdflatex
\documentclass[12pt,a4paper]{article}
\usepackage[a4paper,margin=1in]{geometry}
\usepackage[T1]{fontenc}
\usepackage[utf8]{inputenc}
\usepackage[expansion=false]{microtype}
\usepackage{setspace}
\usepackage{parskip}
\usepackage{titlesec}
\usepackage{xcolor}
\usepackage{listings}
\usepackage{fancyhdr}
\usepackage{lastpage}
\usepackage{enumitem}
\usepackage{hyperref}
\usepackage{graphicx}
\usepackage{float}

\onehalfspacing
\setlength{\parindent}{0pt}
\hypersetup{colorlinks=true,linkcolor=black,urlcolor=blue,
 pdftitle={Blockchain Laboratory - Experiment 6},pdfauthor={Nikhil Sunchu}}
\pagestyle{fancy}
\fancyhf{}
\lhead{Blockchain Laboratory}
\rhead{Experiment No. 6}
\cfoot{Page \thepage}
\setlength{\headheight}{15pt}
\titleformat{\section}{\large\bfseries}{\thesection.}{0.5em}{}
\titleformat{\subsection}{\normalsize\bfseries}{\thesubsection.}{0.5em}{}

\definecolor{codebg}{RGB}{248,248,248}
\definecolor{codecomment}{RGB}{80,120,80}
\definecolor{codekeyword}{RGB}{70,70,150}
\definecolor{codestring}{RGB}{150,70,70}
\definecolor{codetype}{RGB}{0,110,110}
\lstdefinestyle{soliditystyle}{
 language=C,
 backgroundcolor=\color{codebg},
 basicstyle=\ttfamily\footnotesize,
 keywordstyle=\color{codekeyword}\bfseries,
 commentstyle=\color{codecomment},
 stringstyle=\color{codestring},
 numbers=left, numberstyle=\tiny\color{gray}, numbersep=8pt,
 frame=single, framerule=0.4pt, breaklines=true,
 breakatwhitespace=false, showstringspaces=false,
 tabsize=4, columns=fullflexible, keepspaces=true,
 morekeywords={pragma,solidity,contract,mapping,address,uint,public,private,internal,external,view,pure,returns,return,function,constructor,if,else,for,while,require,msg,sender},
 keywordstyle=[2]\color{codetype}\bfseries,
 morekeywords=[2]{address,uint,bool,string,bytes}
}
\begin{document}
\begin{titlepage}
\centering
\vspace*{1.5cm}
{\Large\bfseries BLOCKCHAIN LABORATORY\par}
\vspace{1.2cm}
{\large\bfseries Experiment No. 6\par}
\vspace{0.8cm}
{\LARGE\bfseries Creation and Deployment of Smart Contract and Transactions using Solidity and Remix IDE\par}
\vspace{1.5cm}
\begin{tabular}{rl}
\textbf{Name:} & Nikhil Sunchu \\[0.6cm]
\textbf{Roll No.:} & 31 \\[0.6cm]
\textbf{Class/section:} & D20A \\[0.6cm]
\textbf{Date:} & \hrulefill \\[0.6cm]
\end{tabular}
\vfill
{\large Blockchain Laboratory\par}
\end{titlepage}

\section{Aim}
Creation and deployment of a smart contract and execution of transactions using Solidity and Remix IDE.

\section{Theory}
\subsection{1. Smart Contracts and Solidity}
A smart contract is a program deployed on a blockchain that executes predefined operations when its functions are called. Solidity is a high-level, statically typed programming language commonly used to develop smart contracts for Ethereum-compatible blockchain platforms. Remix IDE provides a browser-based environment for writing, compiling, deploying, and testing Solidity contracts.

\subsection{2. State Variables and Constructor}
State variables store data associated with a deployed contract. In this experiment, \texttt{studentName}, \texttt{studentRoll}, and \texttt{transactionCount} store a student's name, roll number, and the number of update-function calls. The constructor executes once during deployment and initializes the student record using the supplied arguments.

\subsection{3. Functions and Visibility}
Functions define the behaviour of a smart contract. The \texttt{public} visibility specifier allows functions to be called externally. The \texttt{updateName()} and \texttt{updateRoll()} functions modify stored state. Public state variables also have automatically generated getter functions, allowing their values to be read through Remix.

\subsection{4. Deployment and Transactions}
Compilation checks the source code and generates contract artifacts, including bytecode and an ABI. Deployment creates a contract instance in the selected blockchain environment. A state-changing function call is submitted as a transaction and updates the contract's state. In Remix VM, transactions are processed on a simulated blockchain and do not require real Ether.

\subsection{5. Calls and Transaction Count}
Reading a public state variable through its getter is a read operation and does not modify contract state. In this contract, each successful call to \texttt{updateName()} or \texttt{updateRoll()} increments \texttt{transactionCount}. The counter therefore demonstrates the number of update-function executions in the experiment; it does not count the deployment or getter calls.

\section{Tools and Environment Used}
\begin{itemize}
    \item Remix IDE (browser-based development environment)
    \item Solidity compiler, version 0.8.34 as shown in the compiler panel
    \item Remix VM with the Osaka VM setting
    \item Solidity source file: \texttt{StudentRecord.sol}
    \item Remix Solidity Compiler and Deploy \& Run Transactions panels
\end{itemize}

\section{Tasks Performed}
\subsection{Task 1 -- Write the Solidity Program}
\begin{itemize}
    \item Created a Solidity source file named \texttt{StudentRecord.sol} in Remix IDE.
    \item Defined the \texttt{StudentRecord} contract with student name, roll number, and transaction count variables.
    \item Used a constructor to initialize the student record during deployment.
    \item Implemented \texttt{updateName()} and \texttt{updateRoll()} to modify the stored values.
    \item Incremented \texttt{transactionCount} after each update.
\end{itemize}

{\renewcommand{\baselinestretch}{0.9}\selectfont
\begin{lstlisting}[style=soliditystyle]
 // SPDX-License-Identifier: MIT
 pragma solidity ^0.8.20;

 contract StudentRecord {
     string public studentName;
     uint256 public studentRoll;
     uint256 public transactionCount;

     constructor(string memory _name, uint256 _roll) {
         studentName = _name;
         studentRoll = _roll;
         transactionCount = 0;
     }

     function updateName(string memory _name) public {
         studentName = _name;
         transactionCount++;
     }

     function updateRoll(uint256 _roll) public {
         studentRoll = _roll;
         transactionCount++;
     }
 }
\end{lstlisting}
}

\subsection{Task 2 -- Compile the Contract}
\begin{itemize}
    \item Opened the Solidity Compiler panel in Remix IDE.
    \item Selected a compatible Solidity 0.8.x compiler.
    \item Compiled \texttt{StudentRecord.sol} and confirmed successful compilation.
    \item The source code and compiler panel are shown in Figure~\ref{fig:compile}.
\end{itemize}

\begin{figure}[H]
\centering
\includegraphics[width=0.96\textwidth,height=0.42\textheight,keepaspectratio]{image.png}
\caption{StudentRecord source code and successful compilation in Remix IDE}
\label{fig:compile}
\end{figure}

\subsection{Task 3 -- Deploy the Contract}
\begin{itemize}
    \item Opened the Deploy \& Run Transactions panel.
    \item Selected Remix VM as the execution environment.
    \item Supplied \texttt{Nikhil Sunchu} and \texttt{31} as constructor arguments.
    \item Deployed the \texttt{StudentRecord} contract and confirmed that it appeared under Deployed Contracts.
    \item The deployment and decoded constructor arguments are shown in Figure~\ref{fig:deploy}.
\end{itemize}

\begin{figure}[H]
\centering
\includegraphics[width=\textwidth,height=0.40\textheight,keepaspectratio]{swappy-20261019-094304.png}
\caption{Deployment of StudentRecord in Remix VM}
\label{fig:deploy}
\end{figure}

\subsection{Task 4 -- Execute State-Changing Transactions}
\begin{itemize}
    \item Entered \texttt{Nikhil Bikshapati Sunchu} as the argument to \texttt{updateName()} and executed the transaction.
    \item Entered \texttt{32} as the argument to \texttt{updateRoll()} and executed the transaction.
    \item Used the Transact buttons to submit both state-changing function calls.
    \item Remix VM processed the calls as transactions on the simulated blockchain.
\end{itemize}

\begin{figure}[H]
\centering
\includegraphics[width=0.82\textwidth,height=0.32\textheight,keepaspectratio]{swappy-20261019-094648.png}
\caption{Execution of updateName() and updateRoll() transactions}
\label{fig:transactions}
\end{figure}

\subsection{Task 5 -- Verify the Updated Values}
\begin{itemize}
    \item Called the \texttt{studentName()} getter and observed \texttt{Nikhil Bikshapati Sunchu}.
    \item Called the \texttt{studentRoll()} getter and observed \texttt{32}.
    \item Called the \texttt{transactionCount()} getter and observed \texttt{2}.
    \item Confirmed that both update functions modified the stored values and incremented the counter.
    \item The returned values are shown in Figure~\ref{fig:verification}.
\end{itemize}


\begin{figure}[H]
\centering
\includegraphics[width=0.31\textwidth,height=0.27\textheight,keepaspectratio]{swappy-20261019-094648.png}
\hfill
\includegraphics[width=0.31\textwidth,height=0.27\textheight,keepaspectratio]{swappy-20261019-094707.png}
\hfill
\includegraphics[width=0.31\textwidth,height=0.27\textheight,keepaspectratio]{swappy-20261019-094725.png}
\caption{Verification of the updated student name, roll number, and transaction count}
\label{fig:verification}
\end{figure}

\section{Result}
The \texttt{StudentRecord} smart contract was successfully written in Solidity, compiled, and deployed using Remix IDE and Remix VM. The stored student name was updated to \texttt{Nikhil Bikshapati Sunchu}, and the roll number was updated to \texttt{32}. The \texttt{transactionCount} getter returned \texttt{2}, confirming that both state-changing update functions were executed successfully.

\section{Conclusion}
The experiment provided practical experience in creating, compiling, and deploying a Solidity smart contract using Remix IDE. Interacting with the deployed \texttt{StudentRecord} contract demonstrated constructor-based initialization, public state variables, state-changing functions, and getter calls. The experiment also illustrated how transactions can update contract state in a simulated blockchain environment without using real cryptocurrency.
\end{document}
